\documentclass[11pt]{article}
\usepackage[T1]{fontenc}
\usepackage[utf8]{inputenc}
\usepackage{amsmath,amssymb,amsthm}
\usepackage[a4paper,margin=1in]{geometry}
\usepackage[hidelinks]{hyperref}
\setlength{\emergencystretch}{3em}

\newtheorem{theorem}{Theorem}
\newtheorem{lemma}[theorem]{Lemma}
\newtheorem{corollary}[theorem]{Corollary}
\theoremstyle{remark}
\newtheorem{remark}[theorem]{Remark}

\newcommand{\N}{\mathbb N}
\newcommand{\vp}{v_p}
\newcommand{\choosef}[2]{\binom{#1}{#2}}
\newcommand{\BoundaryAt}{\mathsf{BoundaryAt}}
\newcommand{\Realized}{\mathsf{Realized}}
\newcommand{\BoundarySafeAt}{\mathsf{BoundarySafeAt}}
\newcommand{\FactorTableau}{\mathsf{FactorTableau}}
\newcommand{\GSys}{\mathcal G}

\title{Erd\H{o}s Problem 700(i):\\
A Compact Synchronized Carry Characterization}
\author{Nanocodex Erd\H{o}s 700 project}
\date{3 August 2026}

\begin{document}
\maketitle

\begin{abstract}
For
\[
 f(n)=\min_{2\leq k\leq \lfloor n/2\rfloor}
       \gcd\!\left(n,\choosef nk\right),
\]
we give an exact structural characterization of the equality $f(n)=n/q$
for every proper prime-power component $q$ of $n$.  The initial reduction
identifies divisibility-minimal carry obstructions.  The main structural step
then compiles simultaneous realization in all relevant prime bases into one
compact finite integer/Boolean system of selectors, prefix products, digits,
and borrows.  The system has no variable or disjunction for each possible
row or multiplier.  Taking $q$ to be the largest prime divisor gives the
maintained formulation of Erd\H{o}s Problem 700(i).  Taking $q$ to be the
greatest exact prime-power component gives the literal formulation in the
1978 Erd\H{o}s--Szekeres paper, with the necessary prime-power exception.
The reductions and both directions of the compiler are checked in Lean;
novelty and priority remain matters for independent review.
\end{abstract}

\section{Definitions and statement}

Call $k$ \emph{admissible for $n$} when
$2\leq k\leq\lfloor n/2\rfloor$.  For a prime $p$, let
$C_p(n,k)$ be the number of carries in the base-$p$ addition
$k+(n-k)=n$.  By Kummer's theorem,
\[
 C_p(n,k)=v_p\!\left(\choosef nk\right).
\]
For an admissible $k$, define its complementary carry weight by
\[
 W_n(k)=\prod_{p\mid n}p^{\max\{v_p(n)-C_p(n,k),0\}}.       \tag{1}
\]

Fix a positive divisor $q$ of $n$.  Define
\[
 \BoundaryAt(n,q,d)
 \quad\Longleftrightarrow\quad
 d\mid n,\quad q<d,\quad
 \frac d r\leq q\ \text{for every prime }r\mid d.          \tag{2}
\]
Thus $d$ is a divisibility-minimal divisor of $n$ lying numerically above
$q$.  Define
\begin{equation}
 \begin{split}
 \Realized(n,d)\quad\Longleftrightarrow\quad
 &\text{there is an }m\geq1\text{ such that }dm\leq\lfloor n/2\rfloor,\\
 &C_r(n,dm)\leq v_r(n)-v_r(d)
   \quad\text{for every prime }r\mid d.
 \end{split}
\tag{3}
\end{equation}
Finally,
\[
 \BoundarySafeAt(n,q)
 \quad\Longleftrightarrow\quad
 \text{no }d\text{ satisfying (2) is realized}.            \tag{4}
\]
All ranges in (2)--(4) are finite, and the same multiplier $m$ must satisfy
every carry constraint in (3).

\begin{theorem}[Prime-power threshold theorem]\label{thm:threshold}
Let $p$ be prime, let $b\geq1$, and put $q=p^b$.  If $q\mid n$ and $q<n$,
then
\[
 \boxed{
 f(n)=\frac nq
 \quad\Longleftrightarrow\quad
 \BoundarySafeAt(n,q).}
\]
\end{theorem}

\section{The exact complementary weight}

Put $G_n(k)=\gcd(n,\choosef nk)$.  Prime by prime,
\[
 v_p(G_n(k))=\min\!\left(v_p(n),C_p(n,k)\right).
\]
Since $\min(a,c)+\max(a-c,0)=a$, (1) gives the exact identity
\[
 \boxed{G_n(k)W_n(k)=n.}                                  \tag{5}
\]
In particular, $W_n(k)$ is a positive divisor of $n$.

The identity
\[
 n\choosef{n-1}{k-1}=k\choosef nk
\]
shows that $n\mid k\choosef nk$.  Taking valuations gives
\[
 v_p(n)-C_p(n,k)\leq v_p(k)
 \quad\text{whenever the left side is positive},
\]
and hence
\[
 \boxed{W_n(k)\mid k.}                                    \tag{6}
\]
This divisibility is what preserves one common index across all primes.

\section{Every proper prime-power threshold is attained}

Write $n=cq$.  Since $q<n$, we have $c\geq2$, so
$2\leq q\leq\lfloor n/2\rfloor$ and $q$ is admissible.  The binomial
identity
\[
 \choosef{cq}{q}=c\choosef{cq-1}{q-1}                    \tag{7}
\]
will give the exact value at this row.  In base $p$, the lowest $b$ digits
of both $cq-1$ and $q-1=p^b-1$ are $p-1$.  Repeated application of Lucas's
theorem therefore gives
\[
 \choosef{cq-1}{q-1}\equiv1\pmod p.                       \tag{8}
\]
If $A=\choosef{cq-1}{q-1}$, then $A$ is coprime to $q$, even when $p\mid c$.
Consequently
\[
 \gcd\!\left(cq,\choosef{cq}{q}\right)
 =\gcd(cq,cA)=c\gcd(q,A)=c=\frac nq.                      \tag{9}
\]
Combining (5) and (9) yields
\[
 W_n(q)=q.                                                  \tag{10}
\]

It follows that
\[
 f(n)=\frac nq
 \quad\Longleftrightarrow\quad
 W_n(k)\leq q\quad\text{for every admissible }k.           \tag{11}
\]
Indeed, (5) reverses the numerical order between $G_n(k)$ and $W_n(k)$,
while (9)--(10) supply equality at an admissible row.

\section{The boundary antichain}

\begin{lemma}\label{lem:boundary}
If $W\mid n$, then
\[
 q<W
 \quad\Longleftrightarrow\quad
 \text{some }d\mid W\text{ satisfies }\BoundaryAt(n,q,d).
\]
\end{lemma}

\begin{proof}
If $W>q$, choose the numerically least divisor $d$ of $W$ with $d>q$.
Then $d\mid n$.  For every prime $r\mid d$, the proper divisor $d/r$ also
divides $W$, so minimality gives $d/r\leq q$.  Thus $d$ satisfies (2).
The reverse implication follows from $q<d\leq W$.
\end{proof}

For $d\mid n$, valuation comparison in (1) gives
\[
 d\mid W_n(k)
 \quad\Longleftrightarrow\quad
 C_r(n,k)\leq v_r(n)-v_r(d)
 \quad\text{for every prime }r\mid d.                     \tag{12}
\]
If the left side holds, (6) gives $d\mid k$, say $k=dm$ with $m\geq1$.
The admissible endpoint and (12) say exactly that $d$ is realized.  Conversely,
a witness in (3), together with (12), gives $d\mid W_n(dm)$.  Therefore
\begin{equation}
 \begin{split}
 &W_n(k)\leq q\quad\text{for every admissible }k\\
 &\qquad\Longleftrightarrow
 \text{no boundary divisor above }q\text{ is realized}\\
 &\qquad\Longleftrightarrow \BoundarySafeAt(n,q).
 \end{split}
\tag{13}
\end{equation}
Equations (11) and (13) prove Theorem~\ref{thm:threshold}.

\section{The compact synchronized compiler}

The preceding theorem is already an exact characterization, but its semantic
definition quantifies over a shared multiplier.  We now describe the stronger
form that removes any need to list those multipliers in the statement.

Fix an ordered exact prime factorization
\[
 n=\prod_{i=1}^{r}p_i^{a_i}
\]
and a baseline $B$ with $2\leq B\leq n$.  A semantic factor tableau selects
exponents $0\leq e_i\leq a_i$ and one positive integer $T$.  Put
\[
 d=\prod_{i=1}^{r}p_i^{e_i},\qquad K=dT.
\]
The tableau is feasible when
\[
 B<d,\qquad K\leq\lfloor n/2\rfloor,
\]
and, for every active index $i$ with $e_i>0$,
\[
 \frac{d}{p_i}\leq B,
 \qquad
 C_{p_i}(n,K)\leq a_i-e_i.                              \tag{14}
\]
Thus $d$ is a boundary divisor and the one literal row $K=dT$ realizes it.
Conversely every realized boundary divisor determines exactly such an
exponent vector and shared $T$.  Hence
\[
 \FactorTableau(n,B)\text{ is feasible}
 \quad\Longleftrightarrow\quad
 \exists d\,[\BoundaryAt(n,B,d)\wedge\Realized(n,d)].   \tag{15}
\]

The explicit system $\GSys(F,B)$ is a natural-linear presentation of this
tableau.  It contains the following finite fields.
\begin{enumerate}
\item One-hot Boolean selectors choose each exponent
      $e_i\in\{0,\ldots,a_i\}$.
\item Prefix variables build $d$ and $K=dT$ by multiplying only by the fixed
      coefficients $p_i^e$ selected in the preceding step.
\item For every $i$, digit variables expand the same integer $K$ in base
      $p_i$ through height $h_i=\lfloor\log_{p_i}n\rfloor$.
\item Boolean borrow variables encode subtraction of $K$ from $n$ in each
      base.  Their sum is exactly $C_{p_i}(n,K)$.
\item Boundary and budget rows enforce (14) when $e_i>0$ and become inactive
      through checked big-$M$ bounds when $e_i=0$.
\end{enumerate}
Every multiplication in these rows is by input data or a selected fixed
prime power; there is no product of two free integer variables.  The system
is therefore a finite Presburger, equivalently mixed integer/Boolean,
feasibility instance.

\begin{theorem}[Compiler soundness and completeness]\label{thm:compiler}
For every ordered exact factorization $F$ of $n$ and every
$2\leq B\leq n$,
\[
 \boxed{
 \GSys(F,B)
 \quad\Longleftrightarrow\quad
 \FactorTableau(n,B)\text{ is feasible}.}
\]
Consequently, for composite $n>1$,
\[
 \boxed{
 f(n)=\frac{n}{P(n)}
 \quad\Longleftrightarrow\quad
 \neg\GSys(F,P(n)).}                                    \tag{16}
\]
\end{theorem}

\begin{proof}
From a satisfying assignment, one-hotness recovers the exponent vector.
The active prefix equalities recover $d$ and $K=dT$ exactly.  The digit and
borrow rows recover the genuine base-$p_i$ subtraction, so the budget rows
give the carry inequalities in (14).  This constructs a semantic tableau.

Conversely, from a semantic tableau choose the corresponding one-hot
selectors, use the actual prefix products of $d$ and $K$, take the canonical
base-$p_i$ digits of $K$, and take the genuine borrows in $n-K$.  These values
satisfy every selector, bound, prefix, boundary, digit, borrow, and budget
row.  This proves the displayed equivalence.  Equation (15), followed by
Theorem~\ref{thm:threshold} at $q=P(n)$, gives (16).
\end{proof}

This compilation is materially smaller than enumerating the possible
multipliers.  The indexed selector, digit, borrow, and prefix coordinates are
linear in
\[
 \sum_i(a_i+1)+\sum_i(h_i+1)+r.
\]
Since $r$, $\sum_i a_i$, and every $h_i$ are bounded by the binary length of
$n$, there are $O((\log n)^2)$ indexed coordinates.  The fixed coefficients
have $O(\log n)$ bits, giving an $O((\log n)^3)$ direct sparse description.
This is a compact symbolic characterization, not a claim that the resulting
feasibility problem is decidable in polynomial time and not a claim that the
equality cases admit a short list of factorization families.

\section{The two formulations of Part (i)}

Let $P(n)$ be the largest prime divisor of $n$.

\begin{corollary}[Maintained largest-prime formulation]
For every composite $n>1$ and every ordered exact factorization $F$ of $n$,
\[
 \begin{aligned}
 f(n)=\frac{n}{P(n)}
 &\quad\Longleftrightarrow\quad \BoundarySafeAt(n,P(n))\\
 &\quad\Longleftrightarrow\quad \neg\GSys(F,P(n)).
 \end{aligned}
\]
\end{corollary}

\begin{proof}
Apply Theorem~\ref{thm:threshold} with $q=P(n)=P(n)^1$.  Since $n$ is
composite, $P(n)$ is a proper divisor of $n$.  The second equivalence is
Theorem~\ref{thm:compiler}.
\end{proof}

The 1978 paper instead uses
\[
 Q(n)=\max_{p\mid n}p^{v_p(n)},                            \tag{17}
\]
the greatest exact prime-power component of $n$.

\begin{corollary}[Literal 1978 formulation]
For every composite $n>1$ and every ordered exact factorization $F$ of $n$,
\[
 \begin{aligned}
 f(n)=\frac{n}{Q(n)}
 &\quad\Longleftrightarrow\quad
 n\text{ is not a prime power and }\BoundarySafeAt(n,Q(n))\\
 &\quad\Longleftrightarrow\quad
 n\text{ is not a prime power and }\neg\GSys(F,Q(n)).
 \end{aligned}
\]
\end{corollary}

\begin{proof}
If $n$ is not a prime power, an exact component attaining (17) is a proper
prime-power divisor, so Theorem~\ref{thm:threshold} applies.  If
$n=p^a$ is composite, then $Q(n)=n$.  Lucas's theorem shows that every
interior binomial coefficient is divisible by $p$, while Kummer's theorem at
$k=p^{a-1}$ gives valuation exactly one.  Thus $f(p^a)=p>1=n/Q(n)$, so the
equality is impossible.  The second equivalence follows from
Theorem~\ref{thm:compiler} at baseline $Q(n)$.
\end{proof}

\begin{remark}
The two formulations are genuinely different.  For $n=12$, the gcd values
for $k=2,3,4,5,6$ are $6,4,3,12,12$, so $f(12)=3$.  Here
$P(12)=3$ and $Q(12)=4$: the historical equality $f(12)=12/Q(12)$ holds,
whereas the maintained equality $f(12)=12/P(12)$ does not.
\end{remark}


\section{Formal verification and source trail}

\begin{flushleft}
The parameterized theorem is
\path{Erdos700PartI.f_eq_div_primePow_iff_boundarySafeAt}.
The maintained and historical corollaries are respectively
\path{Erdos700PartI.f_eq_div_iff_boundarySafe} and
\path{Erdos700PartI.erdos_700_i_historical}.  They are exposed in
\href{https://github.com/gakonst/nanocodex-erdos700/blob/main/proof/PartIVerify.lean}
{\texttt{proof/PartIVerify.lean}}.  The repository verifier builds the
promoted roots, rejects placeholders, audits dependencies, and runs finite
regressions.  Lean reports only the standard Mathlib axioms
\texttt{propext}, \texttt{Classical.choice}, and \texttt{Quot.sound}; it
reports no \texttt{sorryAx} dependency.

The semantic bridge is
\path{Erdos700PartI.factorTableauFeasible_iff_exists_boundary_realized},
and the explicit compiler theorem is
\path{Erdos700PartI.ExplicitG.explicitG_iff_factorTableauFeasible}.
Lean checks both directions, including one-hot decoding, prefix products,
canonical digit and borrow construction, inactive-row bounds, and the exact
carry-budget identity.

The primary historical source is P.~Erd\H{o}s and G.~Szekeres,
``Some number theoretic problems on binomial coefficients,''
\emph{Australian Mathematical Society Gazette} 5 (1978), 97--99,
\href{https://www.renyi.hu/~p_erdos/1978-46.pdf}{archival PDF}.
\end{flushleft}

\end{document}
